\section{Saturation of the physical normalization}\label{c6:subheading-09}


We now show that the upper bound of Proposition~\ref{prop:gas-upper}
has the correct logarithmic exponent at two parameters. The physical
partition is exactly one at \(y=i\), which makes its normalization
accessible through a determinant. Saturation there propagates to
\(y=1\) by subharmonicity and positivity. Section~\ref{c6:uniform-analytic}
will use both individual asymptotics to exclude zero limits when
sharpening the error to bounded factors.

\begin{auxproposition}[Saturation at the two physical parameters]
For the normalized homogeneous polynomial \(P_m\) and bulk constant
\(I\) of Section~\ref{c6:gas},
\begin{equation}\label{eq:gas-individual-saturation}
\begin{split}
 \log P_m(1)&=Im^2+\frac38\log m+o(\log m),\\
 \log P_m(i)&=Im^2+\frac5{24}\log m+o(\log m).
\end{split}
\end{equation}
Consequently,
\begin{equation}\label{eq:saturation-nesting}
 \mathcal Z_{m,m}(1)=P_m(1)/P_m(i)=m^{1/6+o(1)}.
\end{equation}
\end{auxproposition}

\begin{proof}
We first prove saturation at \(y=i\) by a positive determinant comparison,
then transfer it to \(y=1\).
\paragraph{A reference determinant at zero loop weight.}

At \(y=i\) we have \(G=2f_N^2\), since only the empty polygon system contributes. In its determinantal expression, replace \(h(d)=(\cot(d+b)+\cot(d-b))/2\) by \(h_0(d)=-\lambda\tan(\lambda d)/2\), \(\lambda=\pi/(2b)=4/3\). Denote the resulting normalized homogeneous value (same prefactors as \(P_m(i)\)) by \(P_0\) in this comparison. Then
\[
 P_0=\left[\lambda^4\sin^8(b)/D(0)\right]^{m^2}=e^{I m^2}.
\]
The reference evaluation uses Schur's Pfaffian identity, recorded for example in \cite[Equation~(1.2)]{IshikawaOkadaTagawaZeng2006}. Indeed the Schur Pfaffian of \((X_j-X_i)/(X_j+X_i)\) equals the product over \(i<j\) (clearing denominators gives an alternating polynomial of precisely Vandermonde degree, with coefficient fixed recursively at \(X_2=-X_1\)). Take \(X_j=e^{2i\lambda v_j}\); the determinant with tan thus gives the squared product of tans up to a unit phase. At coincidence use \(B(u)\tan(\lambda u)\to-\lambda\sin^2 b\) and \(D(0)=C\sin^4 b\). All signs give positive value (real Pfaffian squared).

For the second equality Fourier integration with
\(\widehat V(p)\) from \eqref{eq:gas-bulk-symbol} gives
\[
 I=2\int_0^\infty \frac{\cosh(3x)-1}{(2\cosh x-1)\sinh(4x)}\frac{dx}{x}.
\]
The multiplier of \(dx/x\), with \(q=e^{-x}\), is \(\sum_{j=1}^{23}c_j q^j/(1-q^{24})\), with \(c_{24-j}=c_j\) and list \(c_j\) up to \(j=12\):
\((0,2,2,0,-6,-6,0,8,8,2,-6,-8)\).
Expand the geometric series (nonnegative grouped terms), integrate differences of exponentials, and pair symmetric indices using the sine product. This gives
\[
 e^I=\prod_{j=1}^{11}\sin(\pi j/24)^{-c_j}=\lambda^4\sin^8(b)/D(0),
\]
by the listed exponents and ordinary angle formulas. Indeed \(\sum c_j=0\), so the products of \((n+j/24)(n+1-j/24)\) can be taken relative to \((n+1/2)^2\), giving exactly the sine ratios. Here and below the elementary identities for cotangent and sine products follow for example by integrating \(\cot z/(z^2-w^2)\) around expanding squares avoiding its poles (giving the symmetric partial fractions), then integrating the resulting log derivative.

\paragraph{The confluent Gram comparison.}
We compare the true determinant with the reference one before estimating
their ratio. Coalescing the rapidities by Taylor row and column
differences replaces the two determinant kernels by matrices
\((-1)^a h^{(a+d)}(0)/(a!d!)\), \(0\le a,d<2m\), and similarly
for \(h_0\); the common squared Vandermonde cancels. Indeed these are the leading alternating terms in the row and column variables treated separately. Even-even and odd-odd blocks vanish. By cotangent partial fractions, the positive poles of \(h\) are \(b+k\pi,\pi-b+k\pi\), and of \(h_0\) are \((2k+1)b\), \(k\ge0\), all with residue \(1/2\) as for their reflections. Thus the negative odd derivatives of order \(2l+1\) (that is, \(-h^{(2l+1)}(0)\), etc.) are integrals of \(p^{2l+1}\) against the sum of \(e^{-p\cdot\text{pole}}\) over positive poles on \(p>0\). Consequently
\[
 P_m(i)/P_0=\det\!\left[\Pi_m R(p)\Pi_m\big|_{\operatorname{ran}\Pi_m}\right]^2,
 \qquad R(p)=\frac{(1+e^{-2L|p|})(1-e^{-6L|p|})}{1-e^{-8L|p|}}\ge1
\]
where \(\Pi_m\) projects to even polynomials of degree \(\le2m-2\) on the real line with density proportional to \(p/\sinh(bp)\). Indeed these are Gram determinants from the even-odd blocks, positive and nondegenerate.

In coordinates \(x=bp/\pi\) take the density \(w(x)=x/\sinh(\pi x)\). Orthogonal polynomials of degrees \(n\) are given by the Taylor coefficients \(A_n(x)\) of
\[
 F(t,x)=(1+it)^{-1-i x}(1-it)^{-1+i x}.
\]
Their squared norms are \((n+1)/2\) because \(\int w(x)F(t,x)F(u,x)dx=1/[2(1-tu)^2]\) near zero, using \(\int w e^{d x}dx=\tfrac12\sec^2(d/2)\) by geometric series and the cotangent fractions. They have the parity of their degree. At each fixed real \(x\),
\[
 A_{2n}(x)=(-1)^n 2\Re\left[2^{-1+i x}(2n)^{i x}\,a(x)\right]+o(1),\qquad
 |a(x)|^2=\sinh(\pi x)/(\pi x).
\]
Indeed subtract \(2^{-1+i x}(1+it)^{-1-i x}\) and the conjugated-coefficient term: the remainder is bounded analytic inside the unit disk (near each of \(\pm i\) by first-order cancellation), hence has square summable coefficients. The binomial coefficients of the singular factors use \(\prod_{k\le n}(1+i x/k)=n^{i x}(a(x)+o(1))\) by summable log remainders and the sine product. Thus for \(x\ne0\)
\[
 \frac{w(x)}{\log m}\sum_{n<m}\frac{2 A_{2n}(x)^2}{2n+1}\ \longrightarrow\ \frac{1}{2\pi},
\]
since \(\sum_{n=1}^m n^{-1+2ix}\) is bounded by integral comparison.

The log determinant is at least the compressed trace of \(\log R\), by scalar Jensen on each orthonormal eigenvector. Applying Fatou with the nonnegative log yields
\[
 \liminf_m\frac{\log P_m(i)-I m^2}{\log m}
 \ge \frac1\pi\int_{\mathbb R}\log R(\pi x/b)\,dx
 =\frac5{24}
\]
by expanding the three logs as before. This saturates the upper bound at \(r=i\pi/2\).

\paragraph{Propagation from \(y=i\) to \(y=1\).}
The subharmonic functions
\[
 f_m(r)=\frac{\log|P_m(e^r)|-I m^2}{\log m}-\left(3/8+2\Re(r^2)/(3\pi^2)\right)
\]
are bounded above locally by \(o(1)\), with
\(f_m(i\pi/2)\to0\); hence they converge to 0 in local \(L^1\)
throughout \(|\Im r|<.54\pi\). To see the propagation, apply the area
mean inequality on disks centered at a saturating point. The same
argument works at a limit of saturating sequences, by comparing with
slightly larger disks centered on those sequences. Integral convergence
then supplies saturating sequences in overlapping disks, so a finite
chain of disks reaches every compact subset of the strip. In particular
we may choose \(r_m\to0\) with \(f_m(r_m)\to0\).

It remains to pass from these nearby complex points to the real point
zero. Define
\[
 g_m(x)=\frac{\log P_m(e^x)-Im^2}{\log m},\qquad x_m=\Re r_m.
\]
The positive Laurent coefficients supplied by the polygon representation
make \(g_m\) convex, and the triangle inequality gives
\(g_m(x_m)\ge(\log|P_m(e^{r_m})|-Im^2)/\log m\).
Together with Proposition~\ref{prop:gas-upper}, this yields
\(g_m(x_m)\to3/8\). Fix \(a>0\). If \(0\le x_m<a\), convexity
gives
\[
 g_m(x_m)\le(1-x_m/a)g_m(0)+(x_m/a)g_m(a).
\]
Use the analogous inequality with \(-a\) when \(x_m<0\). The upper
bounds at \(a\) and \(-a\), together with \(x_m\to0\), imply
\(\liminf g_m(0)\ge3/8\). The gas upper bound supplies the reverse
limsup. This proves the first asymptotic in
\eqref{eq:gas-individual-saturation}; the determinant comparison and the
upper bound at \(r=i\pi/2\) proved the second. Their difference is
\eqref{eq:saturation-nesting}.
\end{proof}
